\SourceLeaf{69}
\[
\begin{array}{rcccccc}
\dfrac{1-(-1)^m}{2}&\text{wenn }m&\text{gerade ist}&=0,&\text{wenn}&\text{ungerade}&=1\\[5pt]
\dfrac{1-(-1)^{m+1}}{2}&\text{—}&\text{—}&=1,&\text{—}&\text{—}&=0
\end{array}
\]

\noindent und so lässt sich das Gewünschte stets durch Hinzufügen eines solchen Faktors erreichen.

Die Abhandlung selbst bietet ein anschauliches Beispiel, auf das diese
Schreibweise hätte angewendet werden können. Soll nämlich eine Potenz eines
gegebenen Kosinus durch die Kosinus von Vielfachen des Winkels dargestellt
werden, so gilt bekanntlich in der Reihe für gerade Potenzen eine andere Regel
als für ungerade. Sehen wir also, wie sich dieser Unterschied unter eine einzige allgemeine Formel bringen lässt.

Für gerade Potenzen, z.\,B. $\cos\varphi^{2m}$, gilt, wie zuvor gezeigt wurde, folgende Regel.
Das Anfangsglied der nach den Kosinus gerader Vielfacher fortschreitenden Reihe
ist

\[
\frac{1}{2^{2m-1}}\frac{{}^{2m}\mathfrak{B}^{(m)}}{2};
\]

\noindent und das allgemeine Glied $h^{\text{tus}}$ ist

\[
= \frac{1}{2^{2m-1}}\,{}^{2m}\mathfrak{B}^{(m+h)}\cos2h\varphi.
\]

Für die ungeraden Potenzen, deren Glieder nur ungerade Vielfache enthalten,
gilt folgende Regel. Das Anfangsglied der Reihe für $\cos\varphi^{2m-1}$ lautet

\[
= \frac{1}{2^{2m-2}}\,{}^{2m-1}\mathfrak{B}^{(m)}\cos\varphi,
\]

\noindent und das allgemeine Glied $h^{\text{tus}}$ lautet

\[
\frac{1}{2^{2m-2}}\,{}^{2m-1}\mathfrak{B}^{(m+h)}\cos(2h+1)\varphi.
\]

Diesen Ausdrücken entsprechend ist das Anfangsglied der Reihe für
$\cos\varphi^r$ (wobei $r$ ebenso gut eine gerade wie eine ungerade Zahl sein kann)

\[
\left(1-\frac{1-(-1)^{r+1}}{4}\right)\frac{1}{2^{r-1}} . {}^{r}\mathfrak{B}^{\left(\frac{1}{2}r+\frac{1-(-1)^r}{4}\right)}\cos\frac{1-(-1^r}{2}\varphi.
\]

\noindent Ist nämlich $r=2m$, so ergibt sich

\[
\frac{1}{2^{r-1}} = \frac{1}{2^{2m-1}},
\]

\noindent und für den Index des Binomialkoeffizienten gilt

\[
\tfrac{1}{2}r+\frac{1-(-1)^r}{4} = m,\qquad {}^{r}\mathfrak{B}^{\left(\frac{1}{2}r+\frac{1-(-1)^r}{4}\right)}\cos\frac{1-(-1)^r}{2}\varphi = {}^{2m}\mathfrak{B}^{(m)},
\]

\noindent wie zuvor.

